% F13：本书原创矢量示意；依据所在节的定义、证明或明确模型。
\begin{figure}[htbp]
\centering
\begin{tikzpicture}[x=1cm,y=1cm,>=stealth,every node/.style={font=\small},line width=.65pt]

\fill[AnalysisBlue!10](2.5,1.7)circle(1.5);\fill[white](2.5,1.7)circle(1.15);
\draw[AnalysisBlue,thick](2.5,1.7)circle(1.15);\draw[AnalysisBlue,thick](2.5,1.7)circle(1.5);
\draw[->,orange!80!black,thick](2.5,1.7)--(4.35,1.7)node[right]{\(\nabla u\)};
\draw[dashed](2.5,1.7)--(2.5,2.85)node[midway,left]{\(\sqrt t\)};
\node at(2.5,-.5){\(u(x)=|x|^2\)};
\node[align=left]at(9,2.2){水平集长度：\(2\pi\sqrt t\)\\[6pt]梯度长度：\(2\sqrt t\)\\[6pt]均匀概率密度：\(1/\pi\)};
\node at(9,.25){\(\rho(t)=\dfrac{2\pi\sqrt t}{\pi\cdot2\sqrt t}=1\)};
\node at(9,-.5){\(0<t<1\)，临界原点另作零测处理};

\end{tikzpicture}
\caption[径向水平集与余面积分母]{径向水平集与余面积分母。二维单位圆盘上的均匀分布经 \(|x|^2\) 推前后，在 \((0,1)\) 上密度为一。薄环的面积与层值宽度之间由梯度长度连接。}
\label{b1:fig:visual-f13}
\end{figure}
